\begin{abstract}
Post-quantum security is critical 
in the quantum era. 
Quantum computers, along with quantum algorithms, make the standard cryptography based on RSA or ECDSA over FL or Blockchain vulnerable.
The implementation of post-quantum cryptography (PQC) over such systems is poorly understood as PQC is still in its standardization phase.
In this work, we propose a hybrid approach to employ PQC over blockchain-based FL (BFL), where we combine a stateless signature scheme like Dilithium (or Falcon) with a stateful hash-based signature scheme like the extended Merkle Signature Scheme (XMSS).
We propose a linear-based formulaic approach to device role selection mechanisms based on multiple factors to address the performance aspect.
Our holistic approach of utilizing a verifiable random function (VRF) to assist in the blockchain consensus mechanism shows the practicality of the proposed approaches.
The proposed method and extensive experimental results contribute to enhancing the security and performance aspects of BFL systems.

% \devkeypoint{Required Changes: 
% \begin{itemize}
%     \item − Motivate the solution 
%     \item − Formally demonstrate the benefits of a PQC-PQC approach
%     \item − Extend the security proof 
%     \item − Extend the experimental work for evaluating the scalability and include comparison
%     \item  − Explore the design space (optimal values) 
%     \item − Evaluate the impact of the approach in terms of the communication channel, power consumption
% \end{itemize}
% }

% \devkeypoint{R2 The solution does not seem to be required. A single efficient PQC solution would be enough. The benefit and need for this complex construction should be demonstrated. In addition the overall security should be formalized and analysed.}

% \devkeypoint{New Scientific Contribution??}

% \devkeypoint{Too many Ideas}

%  
% % to signature 
% for DS.
% \gangli{is the approach purpose for "signature for DS?''}
% Implementation of proposed selection mechanism 
% also results in overall better performing
% BFL, FL and Blockchain networks 
% by introducing more predictability and 
% % deterministic 
% dynamic nature 
% which provides better planning and understandability.
% Also, for consensus, 
% we propose the use of VRF to decide the
% winning block that gets appended to the blockchain.

% \shiva{PQC, XMSS, vrf, BFL, FL, VRF are never defined in the abstract}
% \dev{done}
%     Performance and security of distributed systems (DS) including blockchain and
%     Federated Learning (FL) is a very crucial aspect in today's digital world.
%     In post quantum era, security of such DS systems is a major channelge while they have their own intrinsic performance 
%     bottleneck challenges.
%     We need to develop new approach and model to understand and investigate post quantum secure DS and jointly resolve their performance and security challenges.
%   To this end, we propose a framework by employing hybrid signature scheme
%     for transforming the performance and security of the blockchain and FL systems. 
%   We  explore and exploit device selection to improve the performance by considering several factors, including physical characteristics of devices, trustability from stake accumulation, convergence contribution etc.
%   We shed light on the practicality and adaptability of integrating hybrid signature approach over the blockchain and FL by rethinking a new simple approach for device selection. Our design ideas and assumptions are validated extensively with experimental results.

\end{abstract}

% \begin{abstract}
% With the post-quantum era around the corner and about to cause a spillover of classical cryptographic systems, 
% blockchain and distributed systems, relying on cryptography, needs a major rethink from a post-quantum lens to enhance their resilience. 
% Post Quantum Cryptography (PQC) is known to respond well to possible quantum threats. 
% However, the bottleneck challenges are in the PQC implementation into the existing distributed systems (DS). 
% So, to move towards that direction of understanding and exploring PQC implementation into 
% DS, in this paper, we propose a hybrid signature scheme approach using two PQC signature schemes.  Building on the blockchain's Proof of Stake (PoS) mechanism and extending Federated Learning (FL), we propose a new approach to device selection in BFL and have slightly touched on different aspects like rewarding protocol, selection of winning miner, dynamic number of active roles etc. Furthermore, we explore and exploit the benefits of FL and Blockchain using this new quantum HBS bridging.
% \end{abstract}

% \keywords{Post Quantum Cryptography (PQC), Federated Learning (FL), Blockchain, Distributed Systems}%

% or neither can they be 
% implemented straightaway without causing any impact on the system.
% Schemes like Dilithium and Falcon are stateless, but with bigger key size in comparison to stateful hash based signature (HBS) schemes like XMSS.
% But schemes like XMSS however cannot be used to sign multiple messages with same key and thus,
% they cannot be used for whole live time of a device.
% So, to address this fundamental issue regarding both stateful and stateless PQC schemes, 
% we have proposed a hybrid signature approach to use both Dilithium (or Falcon) and XMSS in blockchain
% based Federated Learning (BFL) scenario.
% Perhaps Stateful Hash-based signatures (HBS) schemes are the best schemes for 
% future blockchain, as NIST has already acknowledged two stateful schemes, viz. (XMSS) and (LMS).
% Both XMSS and LMS are based on a one-time signature (OTS) and the Merkle tree,
% where the key management and the tree size can affect the scheme's performance. 
% With relevant insights from these developments, we propose a new approach to 
% address the difficulties related to stateful HBS when used in blockchain.
